\subsection{Penalized equations and metric-Laplacian envelopes}
\label{sec:penalized-kahler-envelopes}
When the background is K\"ahler, a constant is an admissible minorant of every continuous obstacle. An exponential penalty then gives a smooth equation whose solutions approach the envelope. Its trace estimate uses only an upper bound for the obstacle Laplacian. Such bounds survive smooth approximations of a finite minimum, even though the minimum itself can have a singular negative second derivative.

In this subsection let $\kappa$ be any fixed K\"ahler form. It is the background of the envelope, while $\omega$ remains the fixed metric used to measure obstacle Laplacians. Write
\[
 \omega=\frac{\mathrm i}{2\pi}h_{i\bar j}\,
                 \mathrm dz_i\wedge\mathrm d\bar z_j,
 \qquad D_h a=h^{i\bar j}a_{i\bar j}=-\frac{La}{4\pi}.
\]
Thus $h$ is the coefficient metric of $2\pi\omega$, and $D_h$ is an auxiliary complex Laplacian. All norms for this fixed metric and for the metric of $\omega$ differ by fixed constants. Write $\kappa=\mathrm i k_{i\bar j}\,\mathrm dz_i\wedge\mathrm d\bar z_j/(2\pi)$, and choose $\gamma>0$ with $k\geq\gamma h$. For a potential $a$ put $g_{i\bar j}=k_{i\bar j}+a_{i\bar j}$; when $g>0$, write $D_g=g^{i\bar j}\partial_i\partial_{\bar j}$ and $H=\operatorname{tr}_h g=\operatorname{tr}_h k+D_ha$.

\begin{lemma}\label{lem:penalized-equation-trace}
Let $f$ be smooth and real, $\beta>0$, and suppose a smooth real $a$ satisfies
\begin{equation}\label{eq:penalized-kahler-equation}
 (\kappa+\ddc a)^n=\mathrm e^{\beta(a-f)}\kappa^n,
 \qquad \kappa+\ddc a>0.
\end{equation}
Then $\min_Xf\leq a\leq\max_Xf$. For fixed $\beta$, the metrics $g,h$ are uniformly equivalent with constants depending only on $X,h,k,\beta,\|f\|_\infty$ and $\sup_XD_hf$. For all sufficiently large $\beta$, the upper bound $H\leq C$ depends only on $X,h,k,\|f\|_\infty$ and $\sup_XD_hf$.
\end{lemma}
\begin{proof}
At a maximum of $a$, its complex Hessian is nonpositive, so the determinant ratio relative to $k$ in \eqref{eq:penalized-kahler-equation} is at most one and $a\leq f$. At a minimum it is at least one. These two points give the value bounds.

We establish the differential inequality used for the trace. At a chosen point take $h$-normal holomorphic coordinates in which $g$ is diagonal with eigenvalues $\lambda_i>0$. The normal-coordinate construction is in the proof of \cref{lem:kahler-trace-star}. Commuting the complex derivatives of the local K\"ahler potentials and differentiating the determinant twice give
\[
 D_gH=\sum_i(\log\det g)_{i\bar i}
       +\sum_{i,k,l}\frac{|\partial_i g_{k\bar l}|^2}
                                {\lambda_k\lambda_l}+E,
 \qquad |E|\leq C H\operatorname{tr}_g h.
\]
Here $E$ consists of the second derivatives of $h^{-1}$ multiplied by $g$ and $g^{-1}$; the first derivatives of $h$ vanish at the point. Their fixed coefficient bounds give the displayed estimate. The determinant identity follows by differentiating $\partial_i\log\det g=\operatorname{tr}(g^{-1}\partial_i g)$ and using $\partial g^{-1}=-g^{-1}(\partial g)g^{-1}$. Also $H_k=\sum_i\partial_i g_{k\bar i}$, by the K\"ahler symmetry. Weighted Cauchy--Schwarz therefore gives
\[
 |\partial H|_g^2
 \leq H\sum_{k,i}\frac{|\partial_i g_{k\bar i}|^2}
                           {\lambda_k\lambda_i}
 \leq H\sum_{i,k,l}\frac{|\partial_i g_{k\bar l}|^2}
                           {\lambda_k\lambda_l}.
\]
The quadratic term consequently absorbs the negative term in $D_g\log H$. Substitute $\log\det g=\log\det k+\beta(a-f)$. Since $H\operatorname{tr}_g h\geq n^2$, the bounded derivatives of $\log\det k$ divided by $H$ are absorbed as well. We obtain
\begin{equation}\label{eq:penalized-log-trace}
 D_g\log H\geq\beta\left(1-\frac{\operatorname{tr}_h k+D_hf}{H}\right)
                                  -C\operatorname{tr}_g h.
\end{equation}
The constant $C$ depends only on the fixed geometry.

Put $B=\sup_X\operatorname{tr}_h k+\sup_XD_hf>0$. Here $\sup_XD_hf\geq0$ by integration. For $Q=\log H-Aa$, the identity $D_ga=n-\operatorname{tr}_g k$ and $\operatorname{tr}_g k\geq\gamma\operatorname{tr}_g h$ give, at a maximum of $Q$,
\begin{equation}\label{eq:penalized-trace-maximum}
 0\geq\beta\left(1-\frac B H\right)
             +(A\gamma-C)\operatorname{tr}_g h-An.
\end{equation}
For fixed $\beta$ take $A=(C+\beta B/n^2+1)/\gamma$. The inequality $1/H\leq\operatorname{tr}_g h/n^2$ gives $\operatorname{tr}_g h\leq An$ there. Thus every eigenvalue is at least $(An)^{-1}$. The determinant bound $\prod_i\lambda_i\leq(\sup_X\det k/\det h)\mathrm e^{2\beta\|f\|_\infty}$ bounds every eigenvalue above there. The maximum property of $Q$ and the value bounds then bound $H$ everywhere. The lower determinant bound $\prod_i\lambda_i\geq(\inf_X\det k/\det h)\mathrm e^{-2\beta\|f\|_\infty}$, together with this global upper bound, gives the global lower eigenvalue bound.

For the assertion uniform in large $\beta$, take instead $A=(C+1)/\gamma$. If $\beta>2An$ and $H>2B$ at the maximum, the right side of \eqref{eq:penalized-trace-maximum} is strictly positive. Hence $H\leq2B$ at that point, and everywhere
\[
 H\leq2B\,\mathrm e^{A\operatorname{osc}_X f}.
\]
Both $A$ and the threshold $2An$ depend only on the fixed forms.
\end{proof}

\begin{theorem}\label{thm:smooth-penalized-kahler-equation}
For every smooth real $f$ and every $\beta>0$, equation \eqref{eq:penalized-kahler-equation} has a unique smooth real solution $a_\beta$.
\end{theorem}
\begin{proof}
We use the equations with obstacle $tf$, $0\leq t\leq1$. At $t=0$ the solution is zero. We first prove the higher estimates for any smooth solution along this path. The preceding lemma bounds the metrics above and below, uniformly in $t$ for fixed $\beta,f$.

Let $\nabla,\widehat\nabla$ be the connections of $g,h$. In holomorphic coordinates their coefficients are $\Gamma^k_{ij}=g^{k\bar l}\partial_i g_{j\bar l}$ and $\widehat\Gamma^k_{ij}=h^{k\bar l}\partial_i h_{j\bar l}$. Define
\[
 \Psi^k_{ij}=\Gamma^k_{ij}-\widehat\Gamma^k_{ij}
             =g^{k\bar l}\widehat\nabla_i g_{j\bar l},
 \qquad S=|\Psi|_g^2.
\]
The difference is a tensor: the inhomogeneous coordinate-change terms in the two connections cancel. Covariant differentiation adds a connection coefficient on each upper holomorphic index and subtracts one on each lower holomorphic index; the conjugate rule applies to antiholomorphic indices. This rule and the derivative of the inverse matrix give $\nabla g=0$.

Here are the tensor identities needed to bound $S$. Put $R^k{}_{i\bar b j}=-\partial_{\bar b}\Gamma^k_{ij}$, with $\widehat R$ defined for $h$, and let $R_j{}^k$ be the Ricci endomorphism of $g$. Differentiating the connection formula gives
\[
 R_{i\bar j}=-\partial_i\partial_{\bar j}\log\det g,
 \qquad
 \nabla_{\bar b}\Psi^k_{ij}=\widehat R^k{}_{i\bar b j}
                                      -R^k{}_{i\bar b j}.
\]
The K\"ahler symmetries yield
$g^{a\bar b}\nabla_a R^k{}_{i\bar b j}=\nabla_i R_j{}^k$:
in $g$-normal coordinates the equality is commutation of the three derivatives in $-\partial_{\bar b}\partial_i g_{j\bar l}$; the first metric derivatives vanish there. Consequently, for the rough Laplacian $\Delta=g^{a\bar b}\nabla_a\nabla_{\bar b}$,
\begin{equation}\label{eq:elliptic-calabi-tensor}
 \Delta\Psi^k_{ij}
  =g^{a\bar b}\nabla_a\widehat R^k{}_{i\bar b j}
                                      -\nabla_iR_j{}^k.
\end{equation}
Commuting $\nabla_a$ and $\nabla_{\bar b}$ on an upper index contributes $R$ and on a lower index contributes $-R$. Thus the reversed rough Laplacian satisfies $\overline\Delta\Psi=\Delta\Psi+\mathcal R\Psi$, where
\[
 (\mathcal R\Psi)^k_{ij}
   =R_i{}^p\Psi^k_{pj}+R_j{}^p\Psi^k_{ip}
                                  -R_p{}^k\Psi^p_{ij}.
\]
Applying the product rule to the squared norm, using $\nabla g=0$ and this commutator, gives the exact scalar identity
\begin{equation}\label{eq:elliptic-calabi-norm}
 D_gS=|\nabla\Psi|_g^2+|\overline\nabla\Psi|_g^2
       +2\operatorname{Re}\langle\Delta\Psi,\Psi\rangle_g
       +\langle\mathcal R\Psi,\Psi\rangle_g.
\end{equation}
The last inner product is real because the Ricci endomorphism is Hermitian on each tensor slot.

The equation identifies the Ricci tensor as
\[
 R_{i\bar j}=-\partial_i\partial_{\bar j}\log\det k
       -\beta\bigl(g_{i\bar j}-k_{i\bar j}-tf_{i\bar j}\bigr).
\]
It is bounded in $g$ norm. Its covariant derivative is bounded by $C(1+|\Psi|_g)$: $\nabla g=0$, while the $g$ derivative of each fixed smooth tensor $k,\operatorname{Ric}(k),\partial\bar\partial f$ is its $h$ derivative plus the connection difference acting on its indices. The same bound holds for $\nabla\widehat R^k{}_{i\bar b j}$. Equations \eqref{eq:elliptic-calabi-tensor} and \eqref{eq:elliptic-calabi-norm} therefore imply
\[
 D_gS\geq-C(S+1).
\]
The trace calculation before \eqref{eq:penalized-log-trace}, now retaining its positive quadratic term, also gives $D_gH\geq cS-C$ with $c>0$. Indeed all background errors are bounded under metric equivalence, the determinant contribution $D_h\log\det g=D_h\log\det k+\beta(H-\operatorname{tr}_h k-tD_hf)$ is bounded, and its quadratic sum bounds $|\widehat\nabla g|_g^2=S$ below by a positive fixed multiple. Choose $K$ so that $Kc>C+1$. At a maximum of $S+KH$ we get $0\geq S-C'$, so $S$ is bounded there. The boundedness of $H$ bounds $S+KH$, hence $S$, everywhere. All constants in this paragraph may depend on $\beta$ and the smooth norm of $f$, but are uniform in $t$.

This bound controls every first derivative of $g$: its holomorphic derivatives follow from $\widehat\nabla g=g\Psi$, and its antiholomorphic derivatives are their Hermitian conjugates. Thus $H$ has a uniform first derivative bound. The fixed-background equation $D_ha=H-\operatorname{tr}_h k$ and the differentiated elliptic estimates in \cref{lem:compact-elliptic-foundations,prop:compact-metric-laplacian-lp} give $W^{3,p}$ bounds for every finite $p>1$. To see the differentiated estimate explicitly, apply the frozen-coefficient $W^{2,p}$ proof to a localized first derivative of $a$; its equation has right side a first derivative of $H$ and commutators with at most second derivatives of $a$. Those second derivatives already have $L^p$ bounds by the undifferentiated estimate. Morrey's inequality gives a uniform $C^{2,\alpha}$ bound when $p>2n/(1-\alpha)$.

The inverse matrix $g^{-1}$ is uniformly Lipschitz. For a real coordinate derivative $v=\partial_ra$, differentiation of the determinant equation gives
\[
 g^{i\bar j}\partial_i\partial_{\bar j}v
   =\beta(v-t\partial_rf)+\partial_r\log\det k
                           -g^{i\bar j}\partial_rk_{i\bar j}.
\]
Differentiate this equation once more and apply the same localized $W^{2,p}$ estimate to $\partial_sv$. The new coefficient derivative is bounded, and multiplies a third derivative of $a$, already bounded in $L^p$. This gives $W^{4,p}$ bounds. Inductively, assume $W^{k,p}$ bounds for every finite $p$, with $k\geq4$. Morrey bounds the derivatives of $a$ through order $k-1$. Differentiate the determinant equation $k-1$ times and put its single highest-order term on the left. Every remaining nonlinear product involves derivatives of $g$ of positive orders summing to at most $k-1$. Its highest possible order is $k-2$, corresponding to a derivative of $a$ of order $k$; there is at most one such factor, and all other factors are bounded. These products are bounded in $L^p$. The frozen-coefficient estimate, whose chart sizes are fixed by the uniform Lipschitz bound of $g^{-1}$, yields $W^{k+1,p}$. Cutoff commutators contain only previously bounded derivatives. Induction and Morrey therefore give every smooth norm, uniformly in $t$.

We supply the openness argument as well. Use the chart Sobolev spaces of \cref{lem:compact-elliptic-foundations}, and fix an integer $s>n+6$. They satisfy
\[
 \|bc\|_{H^s}\leq C_s\|b\|_{H^s}\|c\|_{H^s}.
\]
In a chart this follows from the Fourier convolution formula: the weight $(1+|\xi|^2)^{s/2}$ is bounded by a fixed multiple of the sum of the corresponding weights at $\eta$ and $\xi-\eta$; Young's $L^2*L^1$ estimate and Cauchy--Schwarz bound the unweighted Fourier $L^1$ norm by the $H^s$ norm, since $s>n$. A finite partition proves the compact assertion. For completeness, weighted Cauchy--Schwarz gives $|u*v(x)|^2\leq\|v\|_1\int|u(x-y)|^2|v(y)|\,\mathrm dy$; integration and Fubini give $\|u*v\|_2\leq\|u\|_2\|v\|_1$.

On the open positive-metric subset of $H^{s+2}$, define
\[
 \mathcal N_t(a)=\log\frac{\det(k+a_{i\bar j})}{\det k}
                                                 -\beta(a-tf).
\]
The matrix series for $\log(I+B)$ and $(I+B)^{-1}$ converge in the preceding Sobolev algebra for sufficiently small $B$. They prove that the derivative at a smooth solution is $\mathcal D=D_g-\beta$, and that the Taylor remainder is bounded by $C\|b\|_{H^{s+2}}^2$ with a locally Lipschitz derivative. The map $\mathcal D:H^{s+2}\to H^s$ is invertible. Indeed the positive quadratic form of $-\mathcal D$ is the sum of the gradient integral and $\beta\int|b|^2$, using the volume of $g$. Riesz representation and the difference-quotient proof in \cref{lem:compact-elliptic-foundations} give its inverse and its $H^s\to H^{s+2}$ bound; the positive zeroth-order term controls constants too.

At a solution $a$ for $t_0$, the map
$b\mapsto b-\mathcal D^{-1}\mathcal N_t(a+b)$
maps a sufficiently small closed ball in $H^{s+2}$ into itself and has Lipschitz constant at most $1/2$, provided $|t-t_0|$ is small. This follows from the remainder bound and $\mathcal N_t(a)=\beta(t-t_0)f$. Its iterates are Cauchy by the geometric-series bound on their successive differences, and their limit solves the equation. Sobolev embedding preserves positivity and gives at least six continuous derivatives. The higher-derivative argument above then gives smoothness; difference quotients justify each differentiated equation if the next weak derivative is not yet known. Thus the set of solvable $t$ is open. The uniform smooth estimates and Arzel\`a--Ascoli make it closed: along a sequence of solutions tending to a parameter, a diagonal subsequence converges in every smooth norm, and the uniform lower eigenvalue bound preserves positivity. A nonempty subset of $[0,1]$ that is both open and closed is all of $[0,1]$.

Finally, at a maximum of the difference of two solutions, the first complex Hessian is at most the second one. Their determinant equations give $\beta(a-b)\leq0$ there. Interchanging them proves uniqueness.
\end{proof}
\index{Monge--Amp\`ere equation!smooth exponential penalty}

\begin{proposition}\label{prop:penalized-envelope-uniform-limit}
For smooth real $f$, the solutions above converge uniformly to $P_\kappa(f)$ as $\beta\to\infty$. More precisely, put $B=\sup_X\operatorname{tr}_h k+\sup_XD_hf$, $d_0=\inf_X\det k/\det h$ and $M=B^n/(n^nd_0)\geq1$. Then
\[
 P_\kappa(f)-\frac{\operatorname{osc}_Xf+n\log\beta}{\beta}
  \leq a_\beta
  \leq P_\kappa(f)+\frac{\log M}{\beta},
 \qquad \beta\geq1.
\]
Moreover $P_\kappa(f)$ has bounded metric Laplacian, with a bound depending only on the fixed geometry, $\|f\|_\infty$ and $\sup_XD_hf$.
\end{proposition}
\begin{proof}
The matrix arithmetic--geometric mean bound gives $\operatorname{tr}_h k\geq n d_0^{1/n}$ everywhere. Since $\sup_XD_hf\geq0$, this proves $M\geq1$. At a maximum of $a_\beta-f$, positivity gives $0<g\leq k+(f_{i\bar j})$. The arithmetic--geometric mean bound gives
$\mathrm e^{\beta(a_\beta-f)}\leq M$ there. Subtracting $\log M/\beta$ from $a_\beta$ therefore gives a competitor for $P_\kappa(f)$, proving the upper bound.

The envelope $v=P_\kappa(f)$ is bounded between $\min f$ and $f$. For $0<t\leq1$ put
\[
 w=(1-t)v+t\min_Xf-\frac{n\log(1/t)}{\beta}.
\]
Positive multilinearity gives $(\kappa+\ddc w)^n\geq t^n\kappa^n$, while $w\leq f-n\log(1/t)/\beta$ gives $\mathrm e^{\beta(w-f)}\leq t^n$. Comparison implies $w\leq a_\beta$: on the set $E=\{a_\beta<w\}$, \cref{prop:global-nonpluripolar-comparison} for these bounded potentials would give
\[
 \int_E\mathrm e^{\beta(w-f)}\kappa^n
 \leq\int_E(\kappa+\ddc w)^n
 \leq\int_E\mathrm e^{\beta(a_\beta-f)}\kappa^n.
\]
The first integrand is strictly larger than the last on $E$, so $E$ has zero volume. The almost-everywhere inequality between the two canonical $\kappa$-psh representatives is their everywhere inequality. Taking $t=1/\beta$ proves the lower bound.

The trace lemma supplies $-\operatorname{tr}_h k\leq D_ha_\beta\leq C-\operatorname{tr}_h k$ for all sufficiently large $\beta$. Uniform convergence passes these inequalities to distributions. \Cref{lem:dominated-measure-density} makes the limiting Laplacian an $L^\infty$ function. Its asserted bound follows, and \cref{prop:compact-metric-laplacian-lp} gives the continuous $C^{1,\alpha}$ representative for every $0<\alpha<1$.
\end{proof}

\begin{theorem}\label{thm:kahler-metric-laplacian-rooftop}
Let $r\geq1$ be an integer and let $f_1,\ldots,f_r$ be continuous real functions on $X$ with bounded distributional Laplacians for the fixed metric $\omega$. Then
\[
 P_\kappa\bigl(\min_{1\leq j\leq r}f_j\bigr)
\]
has bounded metric Laplacian. Its bound depends only on the fixed forms $\omega,\kappa$, the uniform norms and $\omega$-metric-Laplacian bounds of the $f_j$. In particular this envelope is $C^{1,\alpha}$ for every $0<\alpha<1$, and its positive form $\kappa+\ddc P_\kappa(\min_jf_j)$ has bounded coefficients. The case $r=1$ gives the same assertion for a single obstacle.
\end{theorem}
\begin{proof}
We first construct smooth approximations with uniform Laplacian bounds. Each $f_j$ is $W^{2,p}$ for every finite $p$ by \cref{prop:compact-metric-laplacian-lp}; fix $p>2n$. Choose a finite smooth partition $\chi_\nu$ supported strictly inside coordinate charts. On a neighborhood of each support take the coordinate convolution $f_{j,\nu,\varepsilon}$ and set $f_{j,\varepsilon}=\sum_\nu\chi_\nu f_{j,\nu,\varepsilon}$. These converge to $f_j$ in $C^1$, by Morrey and smooth approximation.

In real coordinates write $D_h=\sum a^{kl}(x)\partial_k\partial_l+\sum b^k(x)\partial_k$. Its commutator with convolution is the sum of integrals of $(a^{kl}(x)-a^{kl}(y))\partial_k\partial_lf_j(y)$ and $(b^k(x)-b^k(y))\partial_kf_j(y)$ against the kernel. Smooth coefficient bounds and H\"older give an error at most $C\varepsilon^{1-2n/p}\|f_j\|_{W^{2,p}}$. Therefore on each smaller chart
\[
 D_hf_{j,\nu,\varepsilon}
            \leq\operatorname*{ess\,sup}_XD_hf_j+o(1),
\]
and the analogous lower bound holds. In summing the partition, its commutators applied to $f_j$ cancel since $\sum\chi_\nu=1$. Their remaining terms contain only $f_{j,\nu,\varepsilon}-f_j$ and their first derivatives, tending uniformly to zero. Hence the global smooth approximations have uniform upper and lower $D_h$ bounds and uniform value bounds, depending only on the asserted data.

For $\delta>0$ set
\[
 F_{\varepsilon,\delta}
       =-\delta\log\sum_{j=1}^r
                    \mathrm e^{-f_{j,\varepsilon}/\delta},
 \qquad
 p_j=\frac{\mathrm e^{-f_{j,\varepsilon}/\delta}}
                   {\sum_k\mathrm e^{-f_{k,\varepsilon}/\delta}}.
\]
It differs from $\min_j f_{j,\varepsilon}$ by a value in $[-\delta\log r,0]$. The chain rule gives
\[
 D_hF_{\varepsilon,\delta}
 =\sum_jp_jD_hf_{j,\varepsilon}
 -\frac1\delta\left(\sum_jp_j|\partial f_{j,\varepsilon}|_h^2
                    -\left|\sum_jp_j\partial f_{j,\varepsilon}\right|_h^2\right)
 \leq\max_j\sup_XD_hf_{j,\varepsilon}.
\]
The term in parentheses is nonnegative by Cauchy--Schwarz. Thus these smooth obstacles converge uniformly to $\min_j f_j$ while their values and upper Laplacians stay uniformly bounded. The preceding proposition bounds the metric Laplacian of every $P_\kappa(F_{\varepsilon,\delta})$ by the same constant. Uniform envelope stability in \cref{prop:obstacle-envelope-maximal-member} passes the envelopes to the required limit; their distributional Laplacian inequalities pass too. Finally \cref{prop:compact-metric-laplacian-lp,prop:metric-laplacian-regular-contact-density} give the asserted representative and coefficient bounds.
\end{proof}
\index{rooftop envelope!metric-Laplacian regularity}
